\centersection{Įvadas}
\noterequirement{Supažindinama su darbo specifika, aktualumu, išdėstomi tikslai bei uždaviniai, aptariama dokumento struktūra. Šiame skyriuje apie darbą kalbama abstrakčiai, nederėtų pateikti nuorodų į kitus šaltinius (1 – 2 lapai).}

\textit{Darbo problematika ir aktualumas}
\noterequirement{Apibrėžiama darbo problematika ir aptariamas aktualumas. Šiame poskyryje taip pat nurodoma su darbu susijusi sritis, praktinė darbo reikšmė.}

Šiuolaikinis skaitmeninių sistemų luste (\gls{soc}) ir \glsdisp{microprocessor}{mikrovaldiklių} projektavimas tampa vis sudėtingesnis, todėl \glsdisp{verification}{verifikacijos} procesas užima didžiausią laiko dalį viso produkto kūrimo cikle. Tradiciškai rinkoje dominuoja \textit{SystemVerilog} aparatinės įrangos aprašymo kalba ir \glsdisp{ieee}{IEEE 1800.2 \gls{uvm} standartas}, naudojami kartu su brangiais komerciniais \gls{eda} įrankiais. Nors šis standartas užtikrina aukštą našumą, jis sukuria didelį barjerą pradedantiesiems inžinieriams, reikalauja specifinių žinių bei pririša projektus prie konkrečių įrankių tiekėjų\footnote{UVM iš esmės yra nepriklausomas nuo konkrečių įrankių – pats standartas yra nepriklausomas nuo gamintojų. Tačiau praktikoje perkeliant projektą iš vienos EDA priemonės į kitą beveik neabejotinai teks redaguoti dalį kodo, pavyzdžiui atsitiktinių sekų generavimo įrankiai, nors neįtraukiami į UVM standartą, bet beveik visada yra naudojami, o jų realizacija ir elgsena skiriasi skirtinguose EDA įrankiuose \cite{cadence} \cite{synopsis}.}. 

Pastaraisiais metais atviros prieigos bendruomenėse sparčiai populiarėja \textit{Python} programavimo kalba pagrįsti \glsdisp{verification}{verifikavimo} sprendimai (\textit{Cocotb}, \textit{PyUVM}), kurie leidžia naudoti modernaus objektinio programavimo koncepcijas. Tačiau susiduriama su problema su esamu testuojamo skaitmeninio mikrovaldiklio dizaino -- trūksta vieningo, lengvai perkeliamo karkaso, kuris leistų sklandžiai sujungti atviros prieigos \glsdisp{simulator}{imituokles} (pvz., \textit{Verilator}) ir komercinius \gls{eda} įrankius (pvz., \textit{Cadence Xcelium}) į vieną standartizuotą testavimo eigą, išvengiant operacinių sistemų priklausomybių konfliktų. Todėl tokio universalaus, automatizuoto ir \glsdisp{container}{konteinerizuoto} \glsdisp{verification}{verifikavimo} karkaso sukūrimas yra itin aktualus tiek šio universiteto akademinei, tiek šios srities atviro kodo, bendruomenėms.

\textit{Darbo tikslas ir uždaviniai}
\noterequirement{Suformuluojamas pagrindinis darbo tikslas, kuris išskaidomas į kelis uždavinius (3 – 7 uždaviniai). Darbo tikslas turėtų nusakyti kodėl yra kuriama sistema (kokią problemą norima išspręsti), o uždaviniai – kaip to tikslo bus pasiekta (dažniausiai tai sutampa su baigiamojo darbo struktūra – reikės išanalizuoti konkurentus, suformuluoti reikalavimus, ištestuoti sistemą ir t.t.).  Rašant išvadas rekomenduojama atsižvelgti į darbo pradžioje suformuluotus uždavinius. \\
1. Pirmasis uždavinys \\
2. Antrasis uždavinys \\
3. Trečiasis uždavinys \\
}

Pagrindinis šio baigiamojo darbo tikslas -- suprojektuoti ir realizuoti \gls{risc} architektūra pagrįsto \glsdisp{microprocessor}{mikrovaldiklio} automatizuotą, \glsdisp{container}{konteinerizuotą} \glsdisp{verification}{verifikavimo} \glsdisp{framework}{karkasą}, kuris naudoja \textit{Python} \gls{uvm} implementaciją ir užtikrina sklandų testų vykdymą tiek su atviros prieigos, tiek su komercinėmis \glsdisp{simulator}{imituoklėmis} skirtingose operacinėse sistemose.

Šiam tikslui pasiekti buvo iškelti šie uždaviniai:
\begin{enumerate}
	\item Išanalizuoti rinkoje esančius komercinius \gls{eda} \glsdisp{verification}{verifikacijos} sprendimus ir palyginti juos su atviros prieigos \textit{Python} (\textit{PyUVM}, \textit{Cocotb}) galimybėmis.
	\item Suprojektuoti izoliuotą sistemos architektūrą, naudojant \textit{Docker}/\textit{Podman} \glsdisp{container}{konteinerius} bei \textit{GNU Make} automatizaciją, palaikančią \textit{Verilator} ir \textit{Cadence Xcelium} \glsdisp{simulator}{imituokles}.
	\item Realizuoti \gls{uvm} architektūros komponentus (\gls{uvc} agentus, monitorius, tvarkykles, virtualų sekų valdytoją ir etaloninį modelį) \textit{Python} kalba, pritaikant juos \gls{dut} periferijos (\gls{uart}, \gls{gpio}) testavimui.
	\item Pritaikyti atsitiktinių sekų generavimą, jog būtų galima naudoti šiuos komponentus atliekant testuojamo skaitmeninio \gls{soc} dizaino \gls{crv} verifikavimą.
	\item Ištestuoti karkaso universalumą skirtingose pagrindinėse operacinėse sistemose ir procesorių architektūrose ir skirtinguose \gls{eda} įrankiuose.
	\item Įvertinti sukurto karkaso efektyvumą atliekant \gls{soc} dizaino statinio ir \glsdisp{functional-coverage}{funkcinio padengiamumo} bei elgsenos teisingumo matavimus.
\end{enumerate}

\textit{Darbo struktūra}
\noterequirement{Aptariama dokumento struktūra. Nurodoma kiek ir kokių skyrių dokumente yra ir kokia informacija juose pateikiama. Labai svarbu, kad šiame skyriuje nerašytumėte bendrinių dalykų, kurie tinkami kone bet kuriam darbui – pvz. „analizės dalyje apžvelgiami egzistuojantys sprendimai, testavimo skyriuje aprašomas testavimas“. Šis apibendrinimas turi atspindėti būtent jūsų darbo struktūrą, vadinasi – „apžvelgiami egzistuojantys buhalterinės apskaitos paketai bei karkasai, skirti darbui su dideliais duomenimis“ ar pan.}

Dokumentą sudaro keturi pagrindiniai skyriai. \hyperref[section:analysis]{Analizės} skyriuje apžvelgiama darbo problematika, esami komerciniai sprendimai bei pagrindžiamas technologijų pasirinkimas. \hyperref[section:project]{Projekto} skyriuje detalizuojami funkciniai ir nefunkciniai reikalavimai, aprašoma karkaso UML architektūra, \glsdisp{container}{konteinerių} diegimo schemos bei \gls{uvm} komponentų statinės ir dinaminės struktūros. \hyperref[section:testing]{Testavimo} skyriuje aprašomi karkaso suderinamumo testai skirtingose platformose bei gauti \gls{dut} statinio ir \glsdisp{functional-coverage}{funkcinio padengiamumo} rezultatai. \hyperref[section:documentation]{Dokumentacija naudotojui} skyriuje pateikiamos diegimo, naudojimosi bei administravimo instrukcijos.


\textit{Sistemos apimtis}
\noterequirement{Nors tai nėra tiesiogiai įvadinis sistemos aprašas, būtent čia būtina nurodyti jau realizuotos sistemos apimtį. Matai gali būti įvairūs (pasirinkti tinkamus): \\
• Kodo eilučių skaičius (galima išskaidyti pagal į vidinę (angl. backend) ir išorinę (angl. frontend) puses, pagal programavimo kalbas ir kt.) \\
• Testų eilučių skaičius \\
• Darbo valandos \\
• Komponentų, klasių, modulių kiekis \\
• Teikiamų paslaugų skaičius \\
• Sukurtų 3D modelių kiekis ar kita su projektu siejama metrika \\
}

Sukurtas \glsdisp{verification}{verifikavimo} \glsdisp{framework}{karkasas} pasižymi šiais apimties ir integracijos rodikliais:
\begin{enumerate}
	\item Palaikomos operacinės sistemos. 5 aplinkos (\textit{Windows 11 WSL}, \textit{Kubuntu Linux}, \textit{Debian}, \textit{Rocky Linux}, \textit{macOS ARM M4}).
	\item Palaikomos \glsdisp{simulator}{imituoklės}: 2 (\textit{Verilator v5.044}, \textit{Cadence Xcelium v24.03}).
	\item Išeities kodo apimtis: Apie 4 tūkstančiai eilučių gryno kodo išdėstytų per \textit{Python}, \textit{C}, \textit{SystemVerilog}, \textit{Bash} ir \textit{Makefile} failus (neskaitant \gls{dut} šakninio kodo).
	\item Aprašyta daugiau nei 60 \textit{Python} klasių ir sukurta daugiau nei 100 projekto failų\footnote{Šioms metrikoms gauti buvo naudojami \textit{cloc} kodo eilučių analizės ir \textit{tree} failų peržiūros ir analizės programinės įrangos.}.
	\item Pasiektas virš 50\% bendrasis \gls{dut} statinis ir \glsdisp{functional-coverage}{funkcinis padengiamumas}.
\end{enumerate}